\chapter{SWE model}
\label{cap:swe_model}










\section{Precipitation data}

The precipitation dataset, developed by Department of Enviromental Science and Policy of the University of Milan, was set up by retrieving 

% 1086.795 188.0656

\begin{figure}[H]
    \centering
    \includegraphics[width=\linewidth]{swe_model/annual_prec_clima_1991_2020.png}
    \caption{Annual precipitation climatology for the Italian Hydrological part of the Greater Alpine Arc (IH-GAR), computed over the 1991--2020 
    period at a spatial resolution of 30 arc-seconds. The highest totals occur along the Alpine foothills, where moist air masses undergo orographic 
    lifting.}
    \label{fig:annual_prec_clima_1991_2020}
\end{figure}

\begin{figure}[H]
    \centering
    \includegraphics[width=\linewidth]{swe_model/seasonal_prec_clima_1991_2020.png}
    \caption{Seasonal precipitation climatology for the Italian Hydrological part of the Greater Alpine Arc (IH-GAR), computed over the 1991--2020 
    period at a spatial resolution of 30 arc-seconds. Autumn is the wettest seaon on average over the lowlands and the Alpine foothills, whereas the 
    internal Alps experience enhanced precipitations during summer.}
    \label{fig:seasonal_prec_clima_1991_2020}
\end{figure}

\begin{figure}[H]
    \centering
    \includegraphics[width=\linewidth]{swe_model/prec_anomaly_clima_1991_2020.png}
    \caption{Annual precipitation anomalies [mm] averaged over the Italian Hydrological part of the Greater Alpine Arc (IH-GAR) for the period 
    1951--2023, relative to the 1991--2020 climatology. Blue (red) bars indicate positive (negative) anomalies. The black line is a 21-year running 
    mean: its dashed segments, at the beginning and at the end of the record, are computed over a truncated time window.}
    \label{fig:prec_anomaly_clima_1991_2020}
\end{figure}










\section{Temperature data}










% Giustificare il perchè viene utilizzata un'architettura così semplice, alla base del modello insomma
\section{Model description}

The snow model developed during this thesis uses a raster based approach, in which adjacent pixels are treated as independent, thus neglecting every kind of 
horizontal transport. The main drivers of snowpack dynamics are total precipitation and air temperature, available on a daily basis. 

\subsection{Snow accumulation}

The simplest approach to estimate the amount of snowfall on a given day, while having only information on air temperature and total precipitation, 
consists in comparing the mean daily air temperature with a temperature threshold $T_\tau$ in order to discriminate between precipitation phases. In such 
a framework, the variability of air temperature is completely neglected, thus leading to rough estimates. However, the availability of daily maximum and 
minimum temperature data allows for a more realistic evaluation of the fraction of precipitation that reaches the ground in solid form, owning to the fact 
that the total range of variability is known on a daily basis. In particular, the amount of snowfall $S_f$ is evaluated at each grid point as 
\begin{equation}
    S_{f} =
    \begin{cases}
        P
            & \hspace*{1.5em} T_{\max} < T_{\tau} \\[8pt]
        P \cdot \dfrac{T_{\tau} - T_{\min}}{T_{\max} - T_{\min}}
            & \hspace*{1.5em} T_{\min} < T_{\tau} < T_{\max} \\[8pt]
        0
            & \hspace*{1.5em} T_{\min} > T_{\tau},
    \end{cases}
    \label{eq:phase_partitioning_SWE_model}
\end{equation}
where $T_{\max}$ and $T_{\min}$ are the daily maximum and minimum temperature respectively, and $P$ the total amount of precipitation.

\section{Snowmelt}

\begin{equation}
    \text{DDF}_{i}\,=\,\text{DDF}_{ave}\,+\,\text{DDF}_{amp}\cdot\sin\left(\frac{2\pi i}{366}\right)
    \label{eq:swe_ddf}
\end{equation}

\begin{equation}
    \text{M}_i\,=\,\text{DDF}_i\left\{\left(T_{mean,i}\,-\,T_{th,m}\right)\,+\,\ln\left[1\,+\,\exp\left(-\frac{T_{mean,i}\,-\,T_{th,m}}{m_m}\right)\right]\right\}
    \label{eq:swe_melt}
\end{equation}

\begin{equation}
    \text{SWE}_i\,=\,\text{SWE}_{i-1}\,+\,P_{s,i}\,-\,M_i
    \label{eq:swe_swe}
\end{equation}

% Bisogna parlare delle glaciarized area in qualche modo